\documentclass[11pt]{article}
\usepackage[T1]{fontenc}
\usepackage{lmodern,microtype}
\usepackage{amsmath,amssymb,amsthm,mathtools}
\usepackage[letterpaper,margin=1in]{geometry}
\usepackage{booktabs,array,enumitem}
\usepackage{xcolor}
\usepackage[colorlinks=true,linkcolor=blue!55!black,citecolor=blue!55!black,urlcolor=blue!55!black]{hyperref}
\newcommand{\F}{F}
\newcommand{\Z}{\mathbb Z}
\newcommand{\C}{\mathbb C}
\newcommand{\Nm}{\operatorname N}
\newcommand{\Nrd}{\operatorname{Nrd}}
\newcommand{\Irr}{\operatorname{Irr}}
\newcommand{\tH}{\widetilde H}
\newcommand{\tZ}{\widetilde Z}
\newcommand{\Ocal}{\mathcal O}
\newcommand{\Reg}{\operatorname{Reg}}
\newcommand{\Ptag}{\textnormal{\textbf{[P]}}}
\newcommand{\Etag}{\textnormal{\textbf{[E]}}}
\newcommand{\Ctag}{\textnormal{\textbf{[C]}}}
\theoremstyle{plain}
\newtheorem{theorem}{Theorem}[section]
\newtheorem{lemma}[theorem]{Lemma}
\newtheorem{proposition}[theorem]{Proposition}
\newtheorem{corollary}[theorem]{Corollary}
\theoremstyle{definition}
\newtheorem{example}[theorem]{Example}
\newtheorem{remark}[theorem]{Remark}
\title{Wild-place metaplectic character transfer in degree two:\\
Toral positivity and the nonabelian gluing problem\\[0.4em]
\normalsize Draft 0 --- standalone companion to Part IV}
\author{David Joyner \and ChatGPT (OpenAI) \and Claude (Anthropic)}
\date{November 8, 2026}
\emergencystretch=3em
\begin{document}
\maketitle
\begin{abstract}
We initiate a separate local project on irreducible-character
integrality for a scalar-central odd-degree metaplectic cover of
the quaternion division algebra.  The target is the still-open
wild-place assertion that the unitary division-side transform
of Part III carries irreducible characters to genuine
irreducible characters.  We distinguish this conjecture from a
proved uniform quadratic toral permutation-character theorem:
for every quadratic $E/F$ in odd residue characteristic, and
every odd $n$ with $\mu_n\subset F$, the discriminant ratio
$\Delta(t)/\Delta(t^n)$ is a finite permutation character of
$E^\times/F^\times$.  A complete worked example over
$F=\mathbb Q_3(\zeta_3)$ gives positive genuine torus
multiplicities $(3,2,2,2)$ in the unramified case and
$(3,2,1,1,1,1)$ for either ramified quadratic torus, each
of total degree $27$.  We prove an integrality--positivity
criterion for the entire unitary character matrix, explain why
toral positivity does not prove global nonabelian positivity,
and give a staged program using the explicit cocycle, finite
twisted quotients, Clifford theory, and exact character tests.
No genuine $27$-dimensional irreducible of the full covered
division algebra is constructed in this draft.
\end{abstract}
\tableofcontents
\section{Purpose, scope, and status}\label{sec:intro}
The wild division-side character problem is separated from the
conditional global trace comparison in Part~IV of the larger
series.  The following results are the starting point, with
proofs either included here or explicitly attributed to earlier
papers.  The remaining full-group integrality assertion is a
\emph{conjecture}, not a theorem.
\begin{center}\small
\begin{tabular}{@{}p{0.58\linewidth}p{0.28\linewidth}p{0.08\linewidth}@{}}
\toprule
Statement & Source & Status\\\midrule
Unitary local $L_D$ and identity factor & Part III, Draft 9 & [E]\\
Heisenberg commutator and degree on maximal tori & Part I, Draft 8 & [E]\\
Quadratic toral permutation-character theorem & This paper, Theorem~\ref{thm:main} & [P]\\
Three exact cubic wild field-torus calculations & Section~\ref{sec:examplewild} & [P]\\
All-row positivity--integrality equivalence & Proposition~\ref{prop:positive} & [P]\\
Irreducible-character integrality on $\widetilde{D^\times}$ & Section~\ref{sec:plan} & [C]\\
\bottomrule
\end{tabular}
\end{center}

\subsection*{The running example}
The local field $F=\mathbb Q_3(\zeta_3)$ studied here is the wild place $\mathfrak p_3=(1-\zeta_3)$ of the series' running example $K=\mathbb{Q}(\zeta_3)$, $n=3$, $r=2$, $D=(-1,4-\zeta_3)_K$, which is ramified exactly at $\mathfrak p_3$ and $\mathfrak p_7=(3+\zeta_3)$ (see \texttt{notes/mjl\_running\_example.tex} in the project repository). Since $3+\zeta_3\equiv1\pmod{\mathfrak p_3}$ is a square, $D_{\mathfrak p_3}\cong(-1,\varpi)$, which is exactly the model $E_0\oplus E_0j$ with $j^2=\varpi$ of Example~\ref{ex:ranktwo-data}. Target~A of \S\ref{sec:plan} is therefore the first open case of wild integrality in the running example.

\section{Division-side character transform and its normalization}\label{sec:transform}
Let $F$ be nonarchimedean, $D/F$ a division algebra of prime
degree $r$, $H=D^\times$, and $\tH$ its specified cover.
Set $Z=F^\times$ and $\tZ=p^{-1}(Z)$, equipped with compatible
Haar measures on $H/Z\simeq\tH/\tZ$.  The Part III transform is
\begin{equation}\label{eq:LD}
 (L_D\Phi)(z\,t^*)
 =b_D\,\widetilde\theta(z)
  \frac{\Delta(t)}{\Delta(t^n)}\Phi(t),\qquad
 b_D=|n|_F^{-(r-1)/2},
\end{equation}
for regular $t$, $z\in\tZ$; it vanishes at nongood regular lifts.
Here
\begin{equation}\label{eq:Delta}
 \Delta(t)=\left|\det\bigl(1-\operatorname{Ad}(t)
 \mid\mathfrak h/\mathfrak t\bigr)\right|_F^{1/2}
 =\prod_{i<j}\left|
 \frac{(\lambda_i-\lambda_j)^2}{\lambda_i\lambda_j}
 \right|_F^{1/2}.
\end{equation}
This equals Part I's $|D(t)|^{1/2}|\Nrd(t)|_F^{-(r-1)/2}$.
Since the power map on a field torus modulo scalars has index
$|n|_F^{-(r-1)}$, Weyl integration gives $b_D^2$
equal to this index.  The factor $b_D$ is independent of
\emph{common} rescalings of the two identified quotient Haar
measures.  The ratio of discriminants extends at central elements
to $|n|_F^{-r(r-1)/2}$ (Part III, Lemma 3.3), hence
\begin{equation}\label{eq:cD}
 (L_D\Theta_\rho)(1)=c_D\dim\rho,\qquad
 c_D=b_D|n|_F^{-r(r-1)/2}
 =|n|_F^{-(r^2-1)/2}.
\end{equation}
The identity factor is pointwise and does not depend on Haar
measures.  Part III, Theorem 9.1, separately gives the adjoint
orbital-integral/test-function transfer.  Test functions transform
by $\widetilde\theta^{-1}$, not by $\widetilde\theta$.

\subsection{Integral-matrix criterion}\label{subsec:integral}
Choose $\widetilde\theta$ first and put
$\theta(a)=\widetilde\theta(a^*)$.
Part III, Theorem 4.3 and Corollary 5.1, show that $L_D$ is a
unitary \emph{isomorphism} of invariant $L^2$ spaces and that
\begin{equation}\label{eq:matrixchars}
 L_D\Theta_\rho=\sum_\tau a_{\rho\tau}\Theta_\tau,
 \qquad a_{\rho\tau}=\langle L_D\Theta_\rho,
 \Theta_\tau\rangle,
\end{equation}
with each row and column finite.  The characters on either side
are orthonormal bases in their fixed-central-character spaces.

\begin{proposition}[\Etag\ Integral-lattice criterion, Part III]\label{prop:int}
The transform $L_D$ permutes the irreducible characters with
positive signs if and only if every $a_{\rho\tau}$ in
\eqref{eq:matrixchars} is an integer.
\end{proposition}
\begin{proof}
An integral unit vector in $\ell^2$ is a signed basis vector;
unitarity gives injectivity and surjectivity because the inverse
is the adjoint of a unitary map.  The value at identity in
\eqref{eq:cD} is positive and fixes the sign.
This is Part III, Corollary 5.2, with its converse.
\end{proof}

\begin{remark}[The actual obstruction]
For wild $v$, Part III does not prove
$a_{\rho\tau}\in\Z$.  The tame Adams argument works with a
finite quotient $Q$ and a central subgroup $C$, under the
hypothesis $\gcd(n,|Q/C|)=1$ (Part III, Lemma 7.5).
At a wild place that hypothesis need not hold.  For the
simplest illustration take $Q=C_p$, $C=1$, $n=p$: then the
Adams operation sends every linear character to the trivial
character.  This does not prove a nonintegral coefficient for
$L_D$; it only shows why that particular proof method fails.
\end{remark}



\section{A useful global criterion: positivity of every row}\label{sec:positive}
\begin{proposition}[\Ptag\ Positivity versus integrality]\label{prop:positive}
With the unitary character matrix $(a_{\rho\tau})$ of
\eqref{eq:matrixchars}, the following are equivalent:
\begin{enumerate}[label=(\alph*)]
\item $L_D$ permutes irreducible characters with positive signs;
\item every $a_{\rho\tau}$ is an integer;
\item every $a_{\rho\tau}$ is a nonnegative real number;
\item every $L_D\Theta_\rho$ is positive definite on $\tH$.
\end{enumerate}
The implication from positive definiteness to integrality
requires \emph{every} row; positivity of a single row does not
force integrality or irreducibility for that row.
\end{proposition}
\begin{proof}
By Part~III, each row and column of $A$ has $\ell^2$ norm one.
An integer row is a signed basis vector, and its sign is positive
by the identity value \eqref{eq:cD}.  If all entries are
nonnegative, orthogonal rows have disjoint supports.  Every
column meets exactly one row, and since its norm is one its
nonzero entry is one; row norms then imply the permutation
property.  The equivalence of nonnegative character coefficients
and positive definiteness follows from Peter--Weyl/GNS for a
compact-mod-center group with fixed unitary central character.
One row could instead be $(2^{-1/2},2^{-1/2},0,\ldots)$.
\end{proof}
\begin{remark}
At a finite twisted quotient the GNS condition can be checked
by exact positivity of the kernel $(\phi(x^{-1}y))_{x,y}$.
A positive-semidefinite test for one $\rho$ verifies positivity,
\emph{not} the integral-character claim for that $\rho$.
Computing the actual coefficients or obtaining a single-row
virtual-character theorem remains indispensable unless
positivity is proved uniformly for all rows.
\end{remark}

\section{A general quadratic theorem}\label{sec:general-toral}
Let $F$ be a finite extension of $\mathbb Q_p$ with $p$ odd, residue field of cardinality $q$, and let $n\geq1$ be odd with $\mu_n\subset F$. Let $E/F$ be any quadratic field extension, with nontrivial Galois involution $\sigma$. Put
\[
 A=E^\times/F^\times,\qquad G=E^1=\ker(N_{E/F}:E^\times\to F^\times).
\]
Hilbert's Theorem 90 identifies $A\xrightarrow{\sim}G$ via $t\mapsto u=t/\sigma(t)$. The normalized absolute value on $E$ is $|x|_E=q_E^{-v_E(x)}$, so $|x|_E=|N_{E/F}x|_F$.

For noncentral $t\in E^\times$ let
\[
 R_E(t)=\frac{\Delta(t)}{\Delta(t^n)},\qquad
 \Delta(t)=\bigl|\det(1-\operatorname{Ad}(t)\mid D/E)\bigr|_F^{1/2}.
\]
The ratio extends continuously at scalar elements, with value $|n|_F^{-1}$.

\begin{theorem}[\Ptag\ Quadratic toral permutation positivity]\label{thm:main}
For any $E/F,n$ as above, $R_E$ is the character of a finite-dimensional permutation representation of a finite quotient of $A$. In particular,
\[
 R_E(t)=\sum_{\chi\in\widehat A}m_\chi\chi(t),\qquad m_\chi\in\mathbb Z_{\geq0},
\]
where only finitely many $m_\chi$ are nonzero. This includes every wild case $p\mid n$. At the identity $R_E(1)=|n|_F^{-1}$.
\end{theorem}

\subsection{A radial valuation formula}
Write $n=p^a m$ with $(m,p)=1$, and put
\[
 t_\zeta=v_E(1-\zeta),\qquad
 G_k=\{u\in G:v_E(1-u)\geq k\}\quad(k\geq1),\qquad G_0=G.
\]
If $a=0$ then $R_E\equiv1$, so assume $a\geq1$.

\begin{lemma}\label{lem:valuations}
For $u\in G$ and any $\zeta\in\mu_n$ whose prime-to-$p$ component is nontrivial, $|1-\zeta u|_E=1$. For $\zeta\in\mu_{p^a}\setminus\{1\}$,
\[
 v_E(1-\zeta u)=\min\{v_E(1-u),t_\zeta\}.
\]
Consequently, with $k=v_E(1-u)\in\mathbb Z_{\geq0}\cup\{\infty\}$,
\begin{equation}\label{eq:radial}
 R_E(u)=q_E^{\frac12\sum_{\zeta\in\mu_{p^a}\setminus\{1\}}\min(k,t_\zeta)}.
\end{equation}
\end{lemma}
\begin{proof}
The reduction of $u\in E^1$ lies in a norm-one subgroup. In the unramified case its intersection with $k_F^\times$ is $\{\pm1\}$, and in the ramified case the reduction of $u$ is itself in $\{\pm1\}$. A root of unity with nontrivial prime-to-$p$ component has reduction in $k_F^\times\setminus\{\pm1\}$, since $n$ is odd, proving the first assertion.

For a nontrivial $p$-power root $\zeta$, the ultrametric inequality gives the asserted minimum when the two valuations differ. If both equal $t_\zeta$ and $v_E(1-\zeta u)>t_\zeta$, then $u\equiv\zeta^{-1}\pmod{\mathfrak p_E^{t_\zeta+1}}$. Applying $\sigma$ and using $\sigma(u)=u^{-1}$ gives $u^{-1}\equiv\zeta^{-1}$ to the same precision, hence $\zeta\equiv\zeta^{-1}$, a contradiction to $v_E(\zeta-\zeta^{-1})=t_\zeta$ (as $p$ is odd). Thus equality holds.

For $u=t/\sigma(t)$ we have
$\Delta(t)=|N_{E/F}(1-u)|_F^{1/2}=|1-u|_E^{1/2}$.
Factor $1-u^n=\prod_{\zeta\in\mu_n}(1-\zeta u)$ and use the two valuation statements. The identity follows for noncentral $t$ and extends continuously to $u=1$. At $u=1$, the cyclotomic product $\prod_{\zeta\in\mu_{p^a}\setminus\{1\}}(1-\zeta)=p^a$ gives $R_E(1)=|p^a|_E^{-1/2}=|n|_F^{-1}$.
\end{proof}

Define for $k\geq0$
\[
 M_k=q_E^{\frac12\sum_{\zeta\in\mu_{p^a}\setminus\{1\}}\min(k,t_\zeta)},\qquad M_0=1,
\]
and put $T=\max_{\zeta\ne1}t_\zeta$. All values $M_k$ are positive integers because roots $\zeta,\zeta^{-1}$ occur in pairs of equal valuation. We have $M_k=M_T$ for $k\geq T$.

\subsection{The unramified filtration}
Suppose $E/F$ is unramified. Then $q_E=q^2$ and reduction of the norm-one group has order $q+1$; at every deeper level the trace-zero subgroup of the residue-field additive group has order $q$. Thus the filtration satisfies
\[
 [G:G_1]=q+1,\qquad [G:G_k]=(q+1)q^{k-1}\quad(k\geq1).
\]
Write $s_k=\#\{\zeta\in\mu_{p^a}\setminus\{1\}:t_\zeta\geq k\}$. Each $s_k$ is even, and $M_k/M_{k-1}=q^{s_k}$.
Set
\[
 c_k=\frac{M_k-M_{k-1}}{(q+1)q^{k-1}}\quad(1\leq k\leq T).
\]
For $k=1$, $M_0=1$ and $q+1$ divides $q^{s_1}-1$ because $s_1$ is even. For $k>1$ with $s_k>0$, $M_{k-1}$ is divisible by $q^{k-1}$, since at least two roots contribute at least $k-1$ to its exponent. Again $q+1$ divides $q^{s_k}-1$. Hence $c_k\in\mathbb Z_{\geq0}$ (and $c_k=0$ if $s_k=0$).

By telescoping \eqref{eq:radial},
\begin{equation}\label{eq:unram}
 \boxed{R_E=\mathbf1_G+\sum_{k=1}^T c_k\,[G:G_k]\,\mathbf1_{G_k}.}
\end{equation}
Here $[G:G_k]\,\mathbf1_{G_k}$ is precisely the permutation character of $G$ on $G/G_k$.

\subsection{The ramified filtration}
Suppose $E/F$ is ramified. Then $q_E=q$ and $t_\zeta=2v_F(1-\zeta)$ is even. The norm-one group has filtration
\[
 [G:G_1]=2,\qquad G_{2j}=G_{2j+1},\qquad [G:G_{2j}]=2q^j\quad(j\geq1).
\]
For completeness, write $E=F(\sqrt d)$ with $v_F(d)=1$. Hilbert 90 gives $u=(a+b\sqrt d)/(a-b\sqrt d)$; normalize $a,b$ to be integral and not both divisible by the uniformizer of $F$. If $a$ is a unit, scale it to $1$ and write $c=b/a\in\mathcal O_F$: then the depth of $u-1$ is $1+2v_F(c)$. Moreover, two such parameters $c,c'$ define elements congruent modulo $\mathfrak p_E^{2j}$ precisely when $c\equiv c'\pmod{\mathfrak p_F^j}$, so $[G_1:G_{2j}]=q^j$. If $a$ is not a unit and $b$ is, then $u\equiv-1$ modulo $\mathfrak p_E$. These observations give the displayed indices and the fact that positive depths are odd.

Define
\[
 c_1=\frac{M_1-1}{2},\qquad
 c_{2j}=\frac{M_{2j+1}-M_{2j-1}}{2q^j}\quad(1\leq j\leq T/2).
\]
The first coefficient is an integer because $M_1=q^{(p^a-1)/2}$ is odd. For $j\geq1$ write
\[
 \frac{M_{2j+1}}{M_{2j-1}}=q^{e_j},\qquad
 e_j=\frac12\sum_{\zeta\ne1}\bigl(\min(2j+1,t_\zeta)-\min(2j-1,t_\zeta)\bigr)\in\mathbb Z_{\geq0}.
\]
If $e_j=0$, then $c_{2j}=0$. Otherwise some pair $\{\zeta,\zeta^{-1}\}$ has $t_\zeta\geq2j$, so $M_{2j-1}$ is divisible by $q^{2j-1}$, hence by $q^j$. Moreover $q^{e_j}-1$ is even. Therefore $c_{2j}\in\mathbb Z_{\geq0}$.

As all positive depths are odd, telescoping now gives
\begin{equation}\label{eq:ram}
\boxed{R_E=\mathbf1_G+c_1\,[G:G_1] \mathbf1_{G_1}
 +\sum_{j=1}^{T/2}c_{2j}\,[G:G_{2j}]\,\mathbf1_{G_{2j}}.}
\end{equation}
Again each summand is the character of a permutation representation. Equations \eqref{eq:unram} and \eqref{eq:ram} prove Theorem~\ref{thm:main}.


\section{From Fourier coefficients to covered-torus restrictions}
\begin{lemma}[\Ptag\ Fourier test for restrictions to covered field tori]\label{lem:toral-Fourier}
Let $E/F$ be any maximal field in $D$, with $r$ prime and
$(r,n)=1$, and put $A_E=E^\times/F^\times$.  The function
\[
 R_E(t)=\frac{\Delta(t)}{\Delta(t^n)}
\]
extends continuously at scalars, is a locally constant function
on the compact abelian group $A_E$, and has a finite Fourier expansion
\begin{equation}\label{eq:R-Fourier}
 R_E(t)=\sum_{\chi\in\widehat{A_E}}m_\chi\chi(t),
 \qquad
 m_\chi=\int_{A_E}R_E(t)\overline{\chi(t)}\,dt.
\end{equation}
Here the Haar measure of $A_E$ has mass $1$.
Choose a compatible genuine central character
$\widetilde\theta$ of $\widetilde Z$ and set
$\theta(a)=\widetilde\theta(a^*)$.  For any smooth character
$\alpha:E^\times\to\mathbb C^\times$ with
$\alpha|_{F^\times}=\theta$, and for each $\chi\in\widehat{A_E}$,
let $\eta_{\alpha\chi}$ be the unique genuine irreducible
representation of $\widetilde{E^\times}$ whose central character
satisfies
\[
 \zeta_{\alpha\chi}(z\,t^*)=
 \widetilde\theta(z)\alpha(t)\chi(t)
 \qquad(z\in\widetilde Z,\ t\in E^\times).
\]
Define a torus-level class function $\mathcal L_E\alpha$ to equal
$b_D\widetilde\theta(z)R_E(t)\alpha(t)$ on $z t^*\in C_E$
and to vanish outside $C_E$.  Then, as a genuine class function
on the \emph{covered torus},
\begin{equation}\label{eq:toral-branch}
 \mathcal L_E\alpha
 =\sum_\chi m_\chi\,\Theta_{\eta_{\alpha\chi}}.
\end{equation}
For an ordinary character $\Theta_\rho$ of $D^\times$,
its actual $L_D$ transform restricts to the sum of these
$\mathcal L_E\alpha$ with the multiplicities in
$\Theta_\rho|_{E^\times}$.
Consequently $m_\chi\in\mathbb Z_{\geq0}$ for every $\chi$
is sufficient for the restriction of every transferred ordinary
irreducible character to this covered torus to be a genuine
representation character.  It is not sufficient for the full
$\widetilde{D^\times}$ character to be a representation character.
\end{lemma}
\begin{proof}
The extension of $R_E$ at scalars follows from Part~III,
Lemma~3.3.  Since $E/F$ is a degree-$r$ field extension with
$(r,n)=1$, the $n$th power of a noncentral $t$ stays noncentral.
The compactness of $A_E$ and local constancy imply that $R_E$
factors through a finite quotient, giving \eqref{eq:R-Fourier}.

By Part~I the center of the covered torus is
$C_E=\widetilde Z\cdot\{t^*:t\in E^\times\}$.  The displayed
formula for $\zeta_{\alpha\chi}$ is well-defined: the
power map on $A_E$ is injective, and on scalar overlaps the
compatibility $\alpha(a)=\widetilde\theta(a^*)$ is precisely
what is required.  Such genuine central characters extend over
$C_E$, and Stone--von Neumann gives a unique $\eta_{\alpha\chi}$.
Its character vanishes off $C_E$, while on $C_E$ it equals
$b_D\zeta_{\alpha\chi}$ because its degree is $b_D$.
The defining formula \eqref{eq:LD} for $L_D$ has the identical
vanishing and equals $b_D\widetilde\theta(z)R_E(t)\alpha(t)$
on $z t^*$.  Comparing with \eqref{eq:R-Fourier} proves
\eqref{eq:toral-branch}.  Finally, restricting a finite-dimensional
ordinary representation to $E^\times$ decomposes it into
one-dimensional characters, all restricting to $\theta$ on
$F^\times$; summing the displayed formula proves the last claim.
\end{proof}


\section{The explicit cubic wild example}\label{sec:examplewild}
\label{subsec:examplewild}
The following calculation is more informative than the dimension
obstruction alone.  It verifies a \emph{necessary representation-ring
integrality test on every maximal torus}, including the two ramified
quadratic field types; the full-group integrality statement remains open.

\begin{example}[The field, algebra, and cover]\label{ex:ranktwo-data}
Take
\[
 F=\mathbb Q_3(\zeta_3),\qquad \varpi=\zeta_3-1,\qquad
 3=-\zeta_3^{-1}\varpi^2,\qquad q_F=3.
\]
Thus $v_F(3)=2$ and $|3|_F=3^{-2}=1/9$.  The compatible
covering parameters in [JC] and Part~I are respectively
$c=1$ and $c_v=2\pmod3$.  Let $D/F$ be the unique quaternion
division algebra.  With $E_0=F(i)$, $i^2=-1$, its unramified
quadratic maximal subfield, we may write
\[
 D=E_0\oplus E_0j,\qquad j^2=\varpi,\qquad
 jx=\sigma(x)j\quad(x\in E_0),
\]
where $\sigma$ is the nontrivial automorphism of $E_0/F$.
Part~I, Proposition~3.5, gives, for every quadratic maximal
subfield $E/F$, the covered-torus commutator
\begin{equation}\label{eq:ranktwo-comm}
 [x,y]_{\widetilde{E^\times}}=
 (N_{E/F}x,N_{E/F}y)_{F,3}^{\,2}
 (x,y)_{E,3}^{-1}.
\end{equation}
The center $C_E$ of $\widetilde{E^\times}$ is the preimage of
$F^\times E^{\times3}$, and
\[
 [E^\times:F^\times E^{\times3}]=9,
 \qquad b_D=3,\qquad c_D=27.
\]
By the Heisenberg-type Stone--von Neumann theorem
([JC], Proposition~2.4; Part~I, Proposition~3.5), every genuine
irreducible representation of $\widetilde{E^\times}$ has degree $3$.
This is \emph{not} a degree calculation for representations of
$\widetilde{D^\times}$.
\end{example}

\subsection{The unramified quadratic torus}
Put $E_0=F(i)$ as above.  Since $X^2+1$ is irreducible over
$\mathbb F_3$, $E_0/F$ is unramified and its residue field is
$\mathbb F_9$.  For $t=a+bi$ with $a,b\in F$ the reduced norm
and trace give the elementary identity
\begin{equation}\label{eq:quadratic-R}
 R_{E_0}(a+bi)=
 \frac{|a^2+b^2|_F}{|3a^2-b^2|_F}.
\end{equation}
Multiplying $t$ by an element of $F^\times$, suppose
$a,b\in\mathcal O_F$ and at least one is a unit.
Then $a^2+b^2$ is always a unit.  If $b$ is a unit, so is
$3a^2-b^2$.  If $b=\varpi u$, then
\[
 3a^2-b^2=3(a^2+\zeta_3u^2),
\]
and $a$ is a unit; the reduction of $a^2+\zeta_3u^2$ modulo
$\varpi$ is either $1$ or $2$, never $0$.  Hence
\begin{equation}\label{eq:unram-values}
 R_{E_0}(a+bi)=
 \begin{cases}
 9,&b\in\varpi\mathcal O_F,\\
 1,&b\in\mathcal O_F^\times.
 \end{cases}
\end{equation}
The finite residue quotient is
\[
 B_0=k_{E_0}^\times/k_F^\times
 \simeq C_4.
\]
The condition $b\in\varpi\mathcal O_F$ says exactly that the
image of $t$ in $B_0$ is the identity.  Denote its four characters
by $\chi_0=1,\chi_1,\chi_2,\chi_3$.  The regular character of
$C_4$ equals $4$ at the identity and $0$ elsewhere. Thus
\begin{equation}\label{eq:unram-fourier}
 \boxed{\quad
 R_{E_0}=1+2\,\mathrm{reg}_{C_4}
 =3\chi_0+2\chi_1+2\chi_2+2\chi_3.
 \quad}
\end{equation}
This is a positive integral Fourier decomposition of degree
$3+2+2+2=9$; it is stronger than merely checking that
$R_{E_0}$ takes integer values.

\subsection{Both ramified quadratic tori}
Every ramified quadratic extension of $F$ can be written
$E_d=F(j_d)$, $j_d^2=d=u\varpi$ with $u\in\mathcal O_F^\times$.
There are two isomorphism classes, distinguished by the square
class of $u$.  For $t=a+bj_d$,
\begin{equation}\label{eq:ramified-R}
 R_{E_d}(a+bj_d)
 =\frac{|a^2-db^2|_F}{|3a^2+db^2|_F}.
\end{equation}
Again normalize $a,b$ to be integral and not both divisible
by $\varpi$.  The terms $3a^2$ and $db^2$ have different
valuations whenever both $a$ and $b$ are units.  A direct
valuation calculation therefore gives
\begin{equation}\label{eq:ramified-values}
 R_{E_d}(a+bj_d)=
 \begin{cases}
 9,& a\in\mathcal O_F^\times,
       \quad b\in\varpi\mathcal O_F,\\
 3,& a,b\in\mathcal O_F^\times,\\
 1,& a\in\varpi\mathcal O_F,
       \quad b\in\mathcal O_F^\times.
 \end{cases}
\end{equation}
For example, in the original cyclic algebra $d=\varpi$,
\[
 \begin{array}{c|ccc}
 t&j&1+j&1+\varpi j\\\hline
 R(t)&1&3&9\\
 (L_D1)(t^*)&3&9&27
 \end{array}
\]
for the central-character normalization of Example~\ref{ex:ranktwo-data}.
The three values illustrate why the unramified torus alone is not
an adequate test.

There is a natural finite quotient
\[
 B_d=E_d^\times/(F^\times U_1)\simeq C_2\times C_3,
\]
where $U_1\subset\mathcal O_{E_d}^\times$ consists of
$a+bj_d$ with $a\in\mathcal O_F^\times$ and
$b\in\varpi\mathcal O_F$; equivalently
$U_1=\mathcal O_F^\times(1+\mathfrak p_{E_d}^2)$.  The $C_2$ coordinate
records the parity of the $E_d$-valuation, and on even elements
the $C_3$ coordinate is $\bar c\in\mathbb F_3$ for
$1+cj_d$ with $c\in\mathcal O_F$ read modulo $\varpi$.  The multiplication law
\[
 (1+cj_d)(1+c'j_d)
 =(1+dcc')\left(1+\frac{c+c'}{1+dcc'}j_d\right)
\]
shows that $\bar c$ is additive, and $j_d^2=d\in F^\times$
provides the order-two factor.

Write $\epsilon$ for the nontrivial character of $C_2$
and $\xi$ for a character of order $3$ of $C_3$.
The identity \eqref{eq:ramified-values} says that $R_{E_d}$
has values $9$ at the identity of $B_d$, $3$ at its two
nonidentity elements in $C_3$, and $1$ on its three odd elements.
Consequently
\begin{equation}\label{eq:ram-fourier}
 \boxed{\quad
 \begin{aligned}
 R_{E_d}
 &=1+\mathrm{reg}_{C_2}+\mathrm{reg}_{C_2\times C_3}\\
 &=3\cdot1+2\epsilon
   +\xi+\xi^2+\epsilon\xi+\epsilon\xi^2.
 \end{aligned}\quad}
\end{equation}
Here the regular representations are inflated from the indicated
quotients.  All coefficients are nonnegative integers and sum
to $3+2+1+1+1+1=9$.  This applies to \emph{both} ramified
quadratic isomorphism classes, independently of $u$.

\begin{proposition}[\Ptag\ All quadratic torus restrictions pass
an integral-character test]\label{prop:all-three-tori}
In the wild example $F=\mathbb Q_3(\zeta_3)$, $r=2$, $n=3$,
let $D/F$ be the quaternion division algebra.  The function
$R_E$ is the character of a nine-dimensional permutation
representation of a finite quotient of $E^\times/F^\times$,
for every quadratic maximal subfield $E\subset D$.
More precisely, for unramified $E$ it is
$1\oplus 2\,\mathbb C[C_4]$, and for either ramified $E$
it is $1\oplus\mathbb C[C_2]\oplus\mathbb C[C_2\times C_3]$.
For a suitable genuine central character inducing the trivial
ordinary central character, the restriction of $L_D\mathbf1$
to $\widetilde{E^\times}$ is therefore an \emph{actual genuine
representation character} of total dimension $27$, namely
\[
 \begin{array}{c|c}
 E/F&\text{multiplicities of three-dimensional genuine torus irreducibles}\\\hline
 \text{unramified}&(3,2,2,2)\\
 \text{ramified, either type}&(3,2,1,1,1,1).
 \end{array}
\]
These multiplicities are invariant under the nontrivial
$F$-automorphism of each $E$.
\emph{No irreducibility or even integrality for the full
$\widetilde{D^\times}$ is asserted.}
\end{proposition}
\begin{proof}
A $p$-adic field of odd residue characteristic has four
square classes; hence $F$ has exactly three quadratic field
extensions, one unramified and two ramified.  Every quadratic
field extension embeds in the quaternion division algebra.
Equations \eqref{eq:unram-fourier} and \eqref{eq:ram-fourier}
prove the first assertion.  Lemma~\ref{lem:toral-Fourier},
together with the degree-three torus Stone--von Neumann theorem,
converts those Fourier coefficients into the displayed
multiplicities; their weighted dimensions are $3\cdot9=27$.
For the trivial ordinary character one can choose
$\widetilde\theta$ with $\widetilde\theta(a^*)=1$:
the odd-$n$ scalar star kills $\mu_3\subset F^\times$
(Part~I, Section~3), and a faithful character of
$\mu_3$ extends across the abelian scalar central group
modulo the star image.  The nontrivial Galois automorphism
acts by inversion on the respective finite quotients' character
parameters (with the $C_2$ factor fixed), so the displayed
multiplicities are invariant.
\end{proof}



\section{A research plan toward nonabelian gluing}\label{sec:plan}
This section records a proposed research program, not further theorems.
The immediate target is the following sharply formulated special case.
\begin{description}[style=nextline]
\item[Target A: the trivial ordinary representation]
For $F=\mathbb Q_3(\zeta_3)$, $D$ quaternion and $n=3$, decide
whether $L_D\mathbf1$ is the character of a single $27$-dimensional
genuine irreducible of $\tH$, with the prescribed central character.
\item[Target B: the full local statement]
For every ordinary irreducible $\rho$ with a fixed compatible central
character, prove that all coefficients $a_{\rho\tau}$ are integral,
or exhibit a specific nonintegral coefficient.
\item[Target C: uniform rank-two theorem]
Extend Target B to all odd $n$, all $F$ of odd residual
characteristic with $\mu_n\subset F$, and quaternion $D/F$.
\end{description}

\subsection{Stage 1: make the multiplier and finite level explicit}
Use the embedding $D=E_0\oplus E_0j\hookrightarrow M_2(E_0)$,
\[
 x+yj\longmapsto
 \begin{pmatrix}x&\varpi y\\\sigma(y)&\sigma(x)\end{pmatrix},
\]
to restrict [JC, equation (6.1)] to $D^\times$.  Translate the
kernel convention $c=1$ into Part~I's relabeled split-side
parameter $c_v=2\pmod3$.  Check explicitly that the resulting
commutator restricts to \eqref{eq:ranktwo-comm} on both unramified
and ramified maximal tori.  These comparisons must be carried out
at the level of a \emph{cocycle and genuine central character},
not merely at the level of the quotient group $D^\times/F^\times$.

Let $U_D^m=1+\mathfrak P_D^m$.  Seek an open normal congruence
subgroup and compatible splitting such that the fixed-central-
character class function $L_D\mathbf1$ is constant on the
appropriate twisted cosets.  Prove this factorization at a stated
finite level; testing an arbitrary quotient that does not capture
$L_D\mathbf1$ would be inconclusive.  Retain the genuine $\mu_3$
action in the finite central extension or twisted group algebra.

\subsection{Stage 2: exact character or Gram calculation}
For the resulting finite projective quotient, compute
\[
 a_{\mathbf1,\tau}=\langle L_D\mathbf1,\Theta_\tau\rangle
\]
with exact cyclotomic arithmetic and the normalized
fixed-central-character measure.  Cross-check its restricted
character values against the two tables in
Proposition~\ref{prop:all-three-tori}; these branching rules
are necessary constraints, not a uniqueness theorem.
A single nonintegral coefficient disproves Target A as stated.
If all coefficients at a genuinely capturing level are integers,
unitarity and the positive identity value imply Target A.
Alternatively compute the twisted positive-semidefinite Gram
matrix as a diagnostic; positive semidefiniteness \emph{alone}
does not prove Target A.

\subsection{Stage 3: construct a candidate genuine representation}
Try to induce from a suitable covered compact open subgroup
containing a maximal torus and deep principal units.
Use the degree-three Stone--von Neumann representation of
$\widetilde{E^\times}$ and analyze the action of the normalizer
by Clifford theory.  A candidate of degree 27 must have
restrictions with multiplicities $(3,2,2,2)$ on the unramified
torus and $(3,2,1,1,1,1)$ on each ramified type.
Prove irreducibility by an intertwining criterion and then compare
its character on all regular classes with \eqref{eq:LD}.
Existence of compatible torus branchings alone does not suffice.

\subsection{Stage 4: from one row to all rows}
Use the structure of ordinary representations of $D^\times$
(admissible pairs, compact induction, and their character formulas)
to transport the construction from $\rho=\mathbf1$ to arbitrary
$\rho$.  Potential orbit-method arguments on deep $p$-adic
analytic congruence groups require special caution: because
$p\mid n$, multiplication by $n$ on the Lie lattice is not
an automorphism, so orbit dilation by $n^{-1}$ is \emph{not}
an established equivalence.  A full positivity theorem for all
rows would invoke Proposition~\ref{prop:positive} directly.

\subsection{Proof checkpoints and stopping criteria}
A satisfactory positive solution must supply (i) an explicit
cocycle/central-character normalization; (ii) a finite-level
factorization theorem or an independent representation
construction; (iii) exact full-group character coefficients;
and (iv) a uniform all-$\rho$ argument for Target B.
A negative answer requires a verifiable nonintegral
$a_{\rho\tau}$ at a rigorously justified level.
The quadratic toral permutation theorem is already proved and
should not be used as a surrogate for either outcome.

\section{Interface with the larger Jacquet--Langlands program}
The companion problem is local.  Even success on Target B would
not on its own prove the additional global comparison of ordinary
$D^\times$ and the covered group: that comparison remains
conditional on matching functions at exceptional split places,
complex-place normalization, a global orbital-coefficient
identity, and the support/germ hypotheses stated in Part~IV.
Conversely, the global comparison could be established in a
restricted test range without first settling every local
irreducible-character row.  These projects should be developed
in parallel but not conflated.


\section*{Acknowledgments and research-status note}
The computations stem from the collaborative Part~IV project of
D. Joyner, ChatGPT (OpenAI), and Claude (Anthropic).  Claude's
referee report on Part~IV, Draft~5 suggested the positivity
criterion and a uniform rank-two toral analysis.  The general
filtration argument and this separate Draft~0 are prepared for
independent mathematical refereeing.  Sources labeled [E] are
imported from Parts I--III or [JC]; results labeled [P] are
proved here.  All proposed gluing mechanisms are research
programs, and should not be cited as established theorems.
\begin{thebibliography}{99}
\bibitem{JoynerI} D. Joyner, Claude (Anthropic), and ChatGPT (OpenAI),
\emph{A metaplectic Jacquet--Langlands transfer for division
algebras of prime degree, I}, Draft 8 (2026).
\bibitem{JoynerIII} D. Joyner, ChatGPT (OpenAI), and Claude (Anthropic),
\emph{A metaplectic Jacquet--Langlands transfer for division
algebras of prime degree, III}, Draft 9 (2026).
\bibitem{JoynerIV} D. Joyner, ChatGPT (OpenAI), and Claude (Anthropic),
\emph{A metaplectic Jacquet--Langlands transfer for division
algebras of prime degree, IV: Local power maps, toral positivity,
and conditional global comparison}, concluding Draft 6 (2026).
\bibitem{JC} D. Joyner with Claude (Anthropic),
\emph{On the metaplectic analog of Kazhdan's ``endoscopic'' lifting},
corrected version (2026), \texttt{metaplectic\_kazhdan\_revised12.tex}.
\bibitem{FK} Y. Flicker and D. Kazhdan,
\emph{Metaplectic correspondence}, Publ. Math. IH\'ES
\textbf{64} (1986), 53--110.
\bibitem{Zou24} J. Zou,
\emph{Local metaplectic correspondence and applications},
arXiv:2209.05852v3 (December 2024), especially \S3 and Remark~3.8.
\end{thebibliography}
\end{document}
